\section{Distance certificates and Euclidean registration}
\label{sec:v77-registration}

Let $d(s,t)=|A(s)-B(t)|$ be known on a product of scalar intervals;
$A$ and $B$ are unknown regular strictly convex boundary patches.  We
give a finite certificate fixing their common Euclidean gauge and then
recover their entire smooth parametrizations.  The certificate is only
a sufficient condition; in particular it is not a necessity theorem
for circular or conic patches.

Choose six parameters $s_0,\ldots,s_5$ and three parameters
$t_0,t_1,t_2$, all in the interiors of their intervals.  Put
\begin{equation}\label{eq:v77-cross-difference}
 D_{ij}=d(s_i,t_j)^2,\qquad
 X_{ij}=-\tfrac12(D_{ij}-D_{i0}-D_{0j}+D_{00}),
 \quad 1\le i\le5,\ 1\le j\le2.
\end{equation}
Regard the $i$th row as a vector $x_i\in\R^2$, set $x_0=0$, and
write $y_0=0,y_1=(1,0)^\mathsf T,y_2=(0,1)^\mathsf T$.  Define
\begin{equation}\label{eq:v77-rank-five}
 \mathcal M_i=
 (x_{i1}^2,\ 2x_{i1}x_{i2},\ x_{i2}^2,\ -2x_{i1},\ -2x_{i2}),
 \qquad 1\le i\le5.
\end{equation}

\begin{proposition}[A finite certificate for a whole distance patch]
\label{prop:v77-distance}
If $X$ has rank two and the square matrix $\mathcal M$ is invertible,
the kernel $d$ determines both parametrized boundary patches, up to one
common Euclidean isometry.  Their relative translation, angle and scale
are determined.  On the image of actual patches and on fixed bounded
classes with positive singular-value margins, this inverse is locally
Lipschitz from $C^k$ distance kernels to $C^k$ parametrized patches for
every finite $k\ge0$.
\end{proposition}
\begin{proof}
Write $a_i=A(s_i)$ and $b_j=B(t_j)$.  The cross differences are
\[
 X_{ij}=(a_i-a_0)\cdot(b_j-b_0).
\]
Rank two implies that $C=(b_1-b_0\ \ b_2-b_0)$ is invertible.  In
the translation gauge $a_0=0$, set $P=C^{-\mathsf T}$,
$G=P^\mathsf TP>0$, and $h=P^\mathsf Tb_0$.  Then
\begin{equation}\label{eq:v77-anchor-realization}
 a_i=Px_i,\qquad b_j=P^{-\mathsf T}(h+y_j),\qquad
 D_{i0}-D_{00}=x_i^\mathsf TGx_i-2h^\mathsf Tx_i.
\end{equation}
The last five equations are a linear system with coefficient matrix
$\mathcal M$ for the five entries
$(G_{11},G_{12},G_{22},h_1,h_2)$.  They uniquely determine $G,h$.
For realizable data $G$ is positive definite.  Choosing the positive
square root $P=G^{1/2}$ in \eqref{eq:v77-anchor-realization} gives a
canonical realization.  Every other invertible matrix with
$P^\mathsf TP=G$ is an orthogonal left multiple, so the anchors are
unique up to that same orthogonal transformation and the initial
translation.  Realizable data also satisfy
$D_{00}=h^\mathsf TG^{-1}h$ and all the remaining distance identities;
these identities are not extra observed coordinates.

For an arbitrary $s$, subtract the squared distances to $b_0$ from
those to $b_1,b_2$:
\begin{equation}\label{eq:v77-trilateration}
 2(b_j-b_0)\cdot A(s)
 =|b_j|^2-|b_0|^2-\bigl(d(s,t_j)^2-d(s,t_0)^2\bigr),
 \qquad j=1,2.
\end{equation}
The two coefficient vectors are independent.  This recovers all of
$A(s)$.  Select three of its now known noncollinear anchors and use
the transposed formula to recover all of $B(t)$.  Such anchors exist
because $X$ has rank two.  Thus the conclusion concerns full functions,
not a finite list of values.

For stability, the linear solve, positive square root on a compact
positive-definite set, and inversion of the two trilateration matrices
are locally Lipschitz.  Slice evaluation of $d$ is bounded in $C^k$,
and squared-distance differences use products of bounded functions.
Apply these observations to each of the explicit formulas.  They also
define a local reconstruction for nearby nonmodel kernels, although such
outputs need not satisfy every unused distance equation.
\end{proof}

\begin{lemma}[Availability of the certificate]
\label{lem:v77-nonconic-certificate}
If an open arc of $A$ is not contained in any conic, one can choose the
six parameters there so that \eqref{eq:v77-rank-five} is invertible,
for any three distinct parameters on a strictly convex open arc of $B$.
The choices and the inequalities persist under small $C^0$
perturbations of the finitely many anchor positions.
\end{lemma}
\begin{proof}
On the first arc the six functions
$1,A_1,A_2,A_1^2,A_1A_2,A_2^2$ are linearly independent; otherwise
a nonzero polynomial of degree at most two would vanish there.
Their evaluation vectors therefore span $\R^6$.  Choose six vectors
forming a basis.  Subtract the row at the first point from the other
five, translate that point to the origin, and make the invertible linear
coordinate change $x=C^\mathsf T(a-a_0)$.  The six-point evaluation
matrix is invertible exactly when the five-column matrix
\eqref{eq:v77-rank-five} is invertible.  Strict convexity ensures that
three distinct points on the second arc are noncollinear, hence $C$
is invertible.  Invertibility and positive singular values are open
conditions on these finite matrices.
\end{proof}

\subsection{Registration by overlaps}

Recover each distance patch in its own canonical Euclidean frame.
Whenever two copies overlap on the same lifted obstacle, select three
distinct scalar parameters in the overlap and evaluate their two
reconstructed parametrizations, using the given overlap map.  The
three points are noncollinear: a line meets the boundary of a strictly
convex body at at most two points.  There is exactly one Euclidean
isometry carrying the ordered first triple to the ordered second triple.
For realizable data it carries the entire common arc to that arc.

Choose a spanning tree of the connected registration graph and propagate
these isometries from a root frame.  This places every patch.  Non-tree
overlaps are consistency equations; for realizable data they hold
because the actual table supplies one placement.  This is not an
assumption of uniqueness: uniqueness follows from the finite distance
certificates and the uniquely determined triple matches along the tree.
Changing the root frame changes all patches by the same element of
$\mathrm E(2)$.

In the periodic case, let $R_0$, $R_1$, $R_2$ be the placed versions
of the same labelled interval on a representative obstacle and its
translates by the two specified abstract basis elements.  At any common
scalar parameter $s_*$,
\begin{equation}\label{eq:v77-lattice}
 l_j=R_j(s_*)-R_0(s_*),\quad j=1,2,\qquad
 L=(l_1\ l_2),\qquad G_\Lambda=L^\mathsf TL.
\end{equation}
These differences are constant on the common interval for actual data.
They are independent because the table has rank-two lattice.  Complete
boundary representatives together with these vectors determine every
translated obstacle and hence the whole periodic table.

This proves the uniqueness part of Theorem~\ref{thm:v77-main}.
Equality of scalar laws first gives equality of absolute actions and
of distance patches, even if source weights differ.  Their canonical
reconstructions and all tree registrations agree.  No supplied
curvature signature, spatial origin, lattice Gram form, or analytic
continuation is used.
